\section{An unbounded arithmetic proof of the saturation step}
\label{sec:saturation-semantics}

We prove the equivalence for both charts and for either of the two
fixed centers $C_g=(A+\lambda)UE$, where $\lambda=\varepsilon\in\{0,1\}$.
Write $K_{\rm pos}=K_+$ and $G=D_g=ED_f$.  The other notation is that of the complete
equation above.  Throughout this section,
all ten independent coordinates are strictly positive integers.  A
``root'' denotes a positive real square root only inside the proof;
the polynomial evaluator has no root operation.

The inherited input is the accepted ten-positive zero-offset
construction, together with its coherent native/controller/refutation
calibration.  Its proof imports the universal Pell facts and the
source61 lemmas named \texttt{lem:rank}, \texttt{lem:index}, and
\texttt{lem:auxcomplete}.  We state their interfaces at the points
where they are used.  The old strong witness is not retained literally:
the proof will recover $I=U\sqrt E$ and, in the converse direction,
choose a new strong Pell period.  The complete original port tuple
and the complete BFS tuple are separate allowed calibrations;
individual ports are never interchanged.

\subsection{Two roots recover the three required Pell conditions}

\begin{lemma}[Root recovery by a coprime product]
\label{lem:saturation-root-recovery}
At every positive integer zero of the displayed polynomial, there
are positive integers $\tau,f,y$ such that
\begin{equation}
 E=\tau^2,\qquad D_f=f^2,\qquad
 QV^2-(Q-1)y^2=1,
 \label{eq:sat-three-identities}
\end{equation}
and
\begin{equation}
 \frac ZM=\varepsilon_g f\tau+\varepsilon_y y,
 \qquad \varepsilon_g,\varepsilon_y\in\{-1,1\}.
 \label{eq:sat-signed-factor}
\end{equation}
No sign restriction on $Z$ or $V$ is needed for this conclusion.
\end{lemma}

\begin{proof}
Before using the equation, $E>1$, $D_f>c+1$, $Q>1$, $M>0$, and
$\delta=Q-1>0$.  Since $V=D_fb-c$ and $b$ is an integer, either
$b=0$ and $V=-c$, or $|V|\ge D_f-c>1$.  In particular $D_y>0$.
Set
\[
 G=ED_f,\qquad Y_* =\frac{D_y}{\delta}
       =\frac{QV^2-1}{Q-1}.
\]
The two positive rational numbers $G,Y_*$ cannot coincide.  Indeed,
$G=Y_*$ would imply
\[
 Q(V^2-G)=1-G.
\]
Both $Q$ and $G$ are divisible by $E$, so this would give $E\mid1$.
This simple obstruction replaces a radical size or sign separation.

At a zero put $z=Z/M$.  The equation is exactly
\[
 (z^2-G-Y_*)^2=4GY_*.
\]
If $z=0$, then $G=Y_*$, which was excluded.  For $z\ne0$,
\[
 \left(\frac{z^2+G-Y_*}{2z}\right)^2=G.
\]
The complementary rational number has square $Y_*$.  Thus both
positive roots are rational.  Since $G$ is an integer, it is an
integer square.  Moreover
\[
 D_f=1+\Delta U^2E\quad\Longrightarrow\quad\gcd(E,D_f)=1.
\]
A square product of positive coprime integers has square factors.
Write $E=\tau^2$, $D_f=f^2$, and $g=f\tau=\sqrt G$.

For clarity, rational integrality of the remaining adjacent-Pell
root has the following elementary descent.  If $y=a/m>0$ in lowest
terms satisfies $Qv^2-(Q-1)y^2=1$ with integer $v>0$, then
$m^2\mid Q-1$.  For $v>1$ use
\[
 v'=(2Q-1)v-2(Q-1)y,
 \qquad y'=(2Q-1)y-2Qv.
\]
The first coordinate is integral, the same norm is preserved, and
$v<y<Qv/(Q-1)$ gives $-v<v'<v$.  The new first coordinate cannot
vanish.  After taking absolute values, its positive integer size
has decreased.  The reduced denominator of $y'$ is still $m$:
$2Q-1\equiv1\pmod m$ and $a$ is coprime to $m$.  Descent reaches
$v=1$, where $y=1$, so $m=1$.  Apply this argument initially to
$v=|V|$.  The signed factorization now gives
\eqref{eq:sat-signed-factor}.
\end{proof}

\subsection{The old strong witness and the early inequalities}

Define only inside the proof
\begin{equation}
 I=U\tau.
 \label{eq:sat-proof-I}
\end{equation}
Then $I$ is a positive integer, $c^2\mid I$, and
\begin{equation}
 f^2-\Delta I^2=1,\qquad Q=(\Delta I)^2,
 \qquad V=f^2(cT'-R)-c.
 \label{eq:sat-restored-kernels}
\end{equation}
These are exactly the normalized strong and auxiliary identities
of the accepted zero-offset proof.  Its old strong input could be
$i_{\rm old}=I/c^2=i\tau$; it is not the new literal coordinate $i$.

\begin{lemma}[Folding and size before transport]
\label{lem:sat-early-fold}
At every positive zero there are indices $p,n_0\ge1$, with $p$ odd,
such that
\[
 c=\psi_A(p),\qquad
 k=2\psi_{2V_0+1}(n_0),\qquad n_0\ge2.
\]
Consequently
\begin{equation}
 k\ge8V_0+4,
 \qquad c>kY\ge k>8q^5,
 \qquad K_{\rm pos}>0,
 \qquad H>3X>0.
 \label{eq:sat-early-sizes}
\end{equation}
These conclusions do not use the sign of the transport numerator.
\end{lemma}

\begin{proof}
The inherited strong/auxiliary folding lemma is uniform for $A\ge3$.
Its hypotheses are the strong norm, $c\le I$, $c^2\mid I$, the
adjacent auxiliary norm, and $V\equiv-c\pmod f$.  They follow from
\eqref{eq:sat-three-identities}--\eqref{eq:sat-restored-kernels}.
It gives $c=\psi_A(p)$ for odd $p\ge1$.  The unchanged first norm
gives the stated expression for $k$.  If $n_0=1$, then $k=2$;
positivity forces $\eta=\zeta=1$, and $c=2Y+1$ lies strictly
between $\psi_A(1)=1$ and $\psi_A(2)=2A$.  This is impossible.
Thus $k\ge2\psi_{2V_0+1}(2)=8V_0+4$.
The inequalities in \eqref{eq:sat-early-sizes} now use
$V_0=X\mu^2Y^2$, $X\ge q$, and $\mu=q^2+1$.  In particular
$c>2X$, so $K_{\rm pos}=X(3c-2X)+2>0$; the formula
$H=4\mu Y(X+1)-1$ gives $H>3X$.
\end{proof}

The port bounds of the native mask give, for every positive assignment,
\begin{align}
 R&<q^4+(\alpha+2)q^2,\notag\\
 |R|&<q^4+Fq^3+(\alpha+2)q^2,\notag\\
 \frac{|R|}{\mu^2}&<F+\alpha+1.
 \label{eq:sat-mask-bounds}
\end{align}
For the last inequality, multiplying by $\mu^2=(q^2+1)^2$ leaves
the positive difference
\[
 F(q^4-q^3+2q^2+1)+\alpha(q^4+q^2+1)+1.
\]
The first bound and $c>8q^5$ show that
$b=c(T+\alpha)-R>0$, hence $b\ge1$ and $V=D_fb-c>1$.
Put
\[
 u=y-V>0,\qquad \rho=\frac{e}{\mu^2}.
\]
The adjacent norm and $\Delta D_f=Q+\Delta<2Q-2$ give
\[
 u=\frac{V^2-1}{(Q-1)(y+V)}
   <\frac{V}{2(Q-1)}<\frac b\Delta.
\]
Using $c<2kY$, $\Delta>8Y$, $\Delta>\mu^2$, and
\eqref{eq:sat-mask-bounds}, we obtain
\begin{equation}
 u<\frac{kT'}4+F+\alpha+1,
 \qquad \rho<\frac{k}{\mu^2}+F+\alpha+1.
 \label{eq:sat-general-u-rho}
\end{equation}
If $C=q-F-\alpha\ge1$, then
$q(q-F)-C=q\alpha+(q-1)C>0$ and therefore $R>0$.
In that case the sharper inequalities are
\begin{equation}
 u<\frac{kT'}4,
 \qquad \rho<\frac{k}{\mu^2}.
 \label{eq:sat-sharp-u-rho}
\end{equation}
At this stage $\rho$ may still be rational.  No denominator has yet
been cancelled.

\subsection{Why the product center suffices on both charts}

\begin{lemma}[Upper and sharp centers]
\label{lem:sat-centers}
At every positive zero,
\begin{equation}
 g<C_g,
 \qquad 0<C_g-g<(\lambda+1)UE\le2UE.
 \label{eq:sat-center-loss}
\end{equation}
For the sharp center $C_0=AUE$, the stronger estimate
\begin{equation}
 C_0-g>\frac{UE}{4A}>X
 \label{eq:sat-direct-gap}
\end{equation}
also holds.  In particular, for either center and either root sign,
$\varepsilon_g g-C_g+X/q^4<0$.
\end{lemma}

\begin{proof}
Since $I>1$ and $f^2=A^2I^2-I^2+1$,
$(A-1)I<f<AI$.  Multiplication by $\tau$ proves
\eqref{eq:sat-center-loss}.  For the sharper lower bound use
\[
 C_0^2-G=U^2E^2-E.
\]
Here $U^2E>2$ and $g<C_0$, so
$(C_0-g)>UE/(4A)$.  The next estimate is unconditional on the
positive coordinate domain: $k\ge2$, $c>kY$, $U\ge c^2$, and
$E>V_0^2k^2$ imply
\[
 UE>V_0^2k^4Y^2\ge16V_0^2Y^2.
\]
Also $A=\mu Y(X+1)+1<4\mu XY$.  Hence
\[
 \frac{UE}{4A}>
 \frac{V_0^2Y}{\mu X}=X\mu^3Y^5>X.
\]
In particular the direct estimate uses neither $q\mid X$ nor any
decoded source information.
\end{proof}

\subsection{Transport soundness without a large majorant}

On the ordinary chart write $W=1+(q-1)L$ and
\[
 t=\frac{B_t}{q-1},\qquad
 s=\frac{Y}{q^3}>0,\qquad
 \kappa=\frac{K_{\rm pos}}H>0.
\]
The exact normalized carrier and
\eqref{eq:sat-signed-factor} give
\[
 n=(\varepsilon_g g-C_g)+(\varepsilon_y y-V)
       -\kappa-s+Wt+\rho.
\]
By Lemma~\ref{lem:sat-centers}, every zero satisfies
\begin{equation}
 n<u+Wt+\rho.
 \label{eq:sat-transport-upper}
\end{equation}
Suppose $B_t\le-1$.  If $C\ge1$, then $Wt<-L=-kT'$;
\eqref{eq:sat-sharp-u-rho} gives
\[
 n<\frac{kT'}4+\frac{k}{\mu^2}-kT'<0.
\]
If $C\le0$, then $B_t=(K_0+w)C-F\le-F$ and $Wt<-FL$.
By \eqref{eq:sat-general-u-rho},
\[
 n<\frac{kT'}4+\frac{k}{\mu^2}
      +2F+2\alpha+2-FkT'<0.
\]
For completeness, the negative of the expression on the right is
increasing in $F,T,\alpha\ge1$ when $k\ge8$; its minimum is
$3k/2-k/\mu^2-6>0$.  Both cases contradict $n>0$.
We have proved $B_t\ge0$, and then
$(K_0+w)C\ge F>0$ gives $C>0$.

On the direct chart put $W=1+q(q-1)L$.  The normalized carrier differs by the transport
term $-WB/[q(q-1)]$ and the additional term $-X/q^4$.
Lemma~\ref{lem:sat-centers} absorbs the resulting $+X/q^4$ in the
equation for $n$.  Thus the same upper bound holds with
$Wt$ replaced by $WB/[q(q-1)]$.  If $B\le-1$ and $C\ge1$,
this term is smaller than $-L$.  If $C\le0$, then
$B=(qK_0+X)C-qF\le-qF$, so it is smaller than
$-qFL\le-FL$.  The identical two contradictions give $B\ge0$
and $C>0$.  Neither argument assumes the divisibility $q\mid X$.

\subsection{Licenses for the inherited index and projection proof}

We record the order of implications because it prevents circular
cancellation arguments.  Once transport is nonnegative, the native
mask estimate gives $3q<R<q^4-2$.  The recovered strong witness
satisfies $c^2\mid I$ and $c>2R$.  Since $c=\psi_A(p)>2$ and
$A\ge3$, one also has $c>2p$.

The source61 rank lemma says that $c=\psi_A(p)>2$,
$f^2-\Delta I^2=1$, $f,I>0$, and $c^2\mid I$ imply
\[
 f=\chi_A(m),\quad I=\psi_A(m),\quad
 pc\mid m,\quad c\mid m,\quad m>2p.
\]
Put $S=\Delta I$.  Then $S\ge2$, $c\mid S$, and
$S^2=\Delta(f^2-1)$.  The two required auxiliary congruences are
\[
 V\equiv-c\pmod f,
 \qquad V\equiv-R\pmod c;
\]
the second uses $f^2\equiv1\pmod c$.  The source61 index lemma,
with $0<p,R<c/2$, now gives $p=R$.

The index $R$ is odd.  If $q$ were odd, then $J$ would be even and
the native mask formula would make $R$ even.  Thus $q$ is even
and $J$ is odd.  Only now may one use
\[
 \gcd\bigl(\mu^2,q^3(q-1)H\bigr)=1
 \quad\text{or, on the direct chart,}\quad
 \gcd\bigl(\mu^2,q^4(q-1)H\bigr)=1.
\]
Indeed $H\equiv-1\pmod\mu$, and even $q$ makes
$\gcd(q-1,\mu)=1$.  The signed carrier is integral, so the
unchanged denominator argument gives $\mu^2\mid e$.

All following uses of the zero-offset parent are now licensed:
its first-index argument gives $n_0=(R+1)/2$; its offset-free
ratio estimate gives the same half-binomial interval and size bound;
and its whole-product projection gives $X=2^R$.  More explicitly,
on the ordinary chart the integral carrier yields
$H\mid q^3(q-1)K_{\rm pos}$.  Writing
\[
 K_{\rm quad}=X\bigl(X+2(A-2)c\bigr)-1,
 \qquad K_{\rm pos}=HXc-2K_{\rm quad},
\]
one cancels only $2$ (since $H$ is odd).  The universal projection
therefore yields
\[
 H\mid q^3(q-1)(X-2^R)(2^RX+1).
\]
The entire product has absolute value less than $H$ by the retained
ratio and height bounds; hence $X=2^R$.  One does not cancel an
unverified unit modulo $H$.  The direct parent supplies exactly
the corresponding product argument with its extra carrier factor.

The new center causes no change to these interfaces.  Both $C_g$
and $g=f\tau$ are integers, so the only nonintegral terms remain
the old fractions.  On the direct chart, after the main projection
makes $K_{\rm pos}/H$ integral, the residual relation is
\[
 \frac{Y}{q^3}-\frac{WB}{q(q-1)}-\frac{X}{q^4}\in\mathbb Z.
\]
Its numerator modulo $q$ gives $q\mid X$.  Since $X=2^R$,
$q$ is a power of two; $R>3q$ then gives $q^4\mid X$.
Set $w=X/q$ and use $B=qB_t$ to finish the old source decoding.

The remaining parent steps give $R\equiv3\pmod4$, the exact
half-binomial value of $Y$, and all native source divisibilities.
Its final reconstruction uses a \emph{fresh unscaled} source61
parameter $A_{\rm source}=Y(X+1)+2$.  It neither divides the new
strong witness nor identifies this unscaled parameter with $A$.
The fixed native/controller calibration then produces a finite
syntactic ZFC refutation.  This proves soundness for each new
center and chart.

\subsection{Completeness: choose a divisible Pell period}

Suppose a finite ZFC refutation exists.  The parent's canonical
construction gives the outer and first/main data, with
\[
 E=\tau^2,\quad c=\psi_A(R),\quad
 R\equiv3\pmod4,\quad R>3q,
\]
and canonical integers $t\ge0$, $\rho\ge0$, $s>0$, and
$\kappa>0$.  Choose the strong pair anew rather than reusing the
parent's particular index.

Let $N_0=c^2\tau$.  The determinant-one matrix
\[
 M_A=\begin{pmatrix}A&\Delta\\1&A\end{pmatrix}
\]
has finite positive order in the group of invertible matrices over
$\mathbb Z/N_0\mathbb Z$.  At a positive multiple $m$ of that order,
\[
 I=\psi_A(m)>0,\qquad f=\chi_A(m)>0,
 \qquad N_0\mid I,\qquad f^2-\Delta I^2=1.
\]
Set $i=I/(c^2\tau)>0$.  Then the literal $U=ic^2$ satisfies
$U=I/\tau$, and the new $D_f,Q$ are exactly $f^2,(\Delta I)^2$.
This construction requires divisibility of a Pell ordinate by a
prescribed integer, which a finite period supplies.  It requires
no intersection theorem for two unrelated Pell sequences.

The source61 auxiliary-completeness lemma assumes only $A\ge2$,
$R\ge3$, $R\equiv3\pmod4$, $c=\psi_A(R)$,
$f^2-\Delta I^2=1$, $f,I>0$, and $c\mid I$.  All hold for this
new pair.  In particular the lemma fixes neither the strong index
nor the parity of $I/c^2$.  Its explicit witness, with
$S=\Delta I$, $Q=S^2$, and $R=2v+1$, is
\begin{align*}
 V_*&=\chi_S(R)/S=Q_v(S^2),\qquad y=\psi_S(R),\\
 T_*&=\frac{V_*+c+Rf^2}{cf^2},\qquad
 QV_*^2-(Q-1)y^2=1.
\end{align*}
The lemma makes $T_*$ a positive integer.  Its usual Pell growth
bound gives, since canonical $v\ge2$,
\[
 V_*\ge(4Q-3)^v\ge(3Q)^v\ge9Q^2.
\]
Because $c<S$, $f^2<Q$, and $\alpha<q<S$,
\[
 T_*>\frac{V_*}{cf^2}>\frac{V_*}{S^3}\ge9S>\alpha.
\]
Thus $T=T_*-\alpha$ is positive and the literal auxiliary register
equals $V_*$.  For $u=y-V_*>0$, using $V_*\ge2$ and $y<2V_*$,
\begin{equation}
 u=\frac{V_*^2-1}{(Q-1)(y+V_*)}
   >\frac{V_*}{4Q}\ge\frac94S^2
   =\frac94\Delta^2I^2.
 \label{eq:sat-complete-gap}
\end{equation}

Here $U=ic^2>1$, so $UE=I\tau<I^2$.  Also
$\kappa<c<I$ and $s<Y<c<I$; the first bound follows from
$K_{\rm pos}<3Xc$ and $H>3X$.  Define the integer
\[
 n=(g-C_g)+(y-V_*)-\kappa-s+Wt+\rho.
\]
Its total negative loss is less than
\[
 2UE+\kappa+s<2I^2+2I<4I^2,
\]
whereas \eqref{eq:sat-complete-gap} exceeds $4I^2$, since
$\Delta\ge3$.  Thus $n>0$.  The normalized center equals
$g+y$, and one factor of the two-root norm vanishes.  This is
completeness on the ordinary chart for both centers.  The explicit
auxiliary witness provides the needed growth; no unproved
cofinality assertion has been inserted.

For the direct chart, transfer this newly constructed ordinary zero.
It has $X=2^R>q$, $q^4\mid X$, and $t=B_t/(q-1)\ge0$.
Set
\begin{align*}
 x&=X-q,\\
 W_{\rm ord}&=1+(q-1)L,\qquad
 W_{\rm dir}=1+q(q-1)L,\\
 n_{\rm dir}&=n_{\rm ord}
       +(W_{\rm dir}-W_{\rm ord})t+X/q^4.
\end{align*}
Every coordinate is positive.  All kernels and $C_g$ agree, and
the normalized centers agree.  In fact
$M_{\rm dir}=qM_{\rm ord}$ and $Z_{\rm dir}=qZ_{\rm ord}$,
so the new norm scales by $q^4$ and remains zero.

\begin{theorem}[Semantic equivalence of the four saturated equations]
\label{thm:saturation-equivalence}
For either fixed center and either chart, with any one complete
coherent calibrated port tuple, the constructed polynomial has a
zero in ten strictly positive integers if and only if there is a
finite syntactic ZFC refutation.  Consequently ZFC proves that its
lack of positive integer solutions is equivalent to
$\operatorname{Con}(\mathrm{ZFC})$.
\end{theorem}

\begin{proof}
Soundness and completeness were proved above, using the same
arithmetizable source calibration as the parent.  Negating the
resulting existence equivalence gives the asserted consistency
equivalence.  The argument is ordinary mathematical proof
formalizable in ZFC; it is not an assertion that a new Lean, Rocq,
or encoded internal-ZFC proof artifact has been produced.
\end{proof}
